\documentclass[a4paper, 10pt]{beamer}

% Encoding
\usepackage[T1]{fontenc}
\usepackage[utf8]{inputenc}
\usepackage[english]{babel}

% Margins
% \usepackage[left=2.5cm, right=2.5cm]{geometry}

% Mathematics
\usepackage{amsmath, amssymb, mathtools, amsthm}
\usepackage{dsfont}

% Bibliography (uncomment if needed)
%\usepackage{csquotes}
%\usepackage[backend=biber, sorting=none, style=nature]{biblatex}
%\addbibresource{biblio.bib}

% Personal taste
\usepackage[font=footnotesize]{caption}
\usepackage[parfill]{parskip}


\title{utils}
\author{}

\newcommand{\PP}[0]{\mathbb{P}}
\newcommand{\Var}[1]{\text{Var}\left(#1\right)}
\newcommand{\E}[1]{\mathbb{E}\left[#1\right]}
\newcommand{\CF}[1]{\mathds{1}_{#1}}

\begin{document}

\begin{frame}{Independence}
  Suppose we are dealing with two r.v. $X$ and $Y$
  \begin{equation*}
    X \sim f_X, \quad Y \sim f_Y
  \end{equation*}
  We say that $X$ and $Y$ are \textbf{independent} if and only if
  \begin{equation*}
    f_{X|Y}(x|y) = f_X(x)
  \end{equation*}
  or, equivalently\footnote{by definition $f_{X,Y}(x,y) = f_{X|Y}(x|y) f_Y(y)$}
  \begin{equation*}
    f_{X,Y}(x,y) = f_X(x)f_Y(y)
  \end{equation*}
\end{frame}

\begin{frame}
  \begin{tabular}{c | c}
    X & Y \\
    \hline
    $x_1$ & $y_1$\\
    $x_2$ & $y_2$\\
    ... & ...\\
    $x_n$ & $y_n$\\
  \end{tabular}
  $\quad \longrightarrow \quad s = s(x_1,...x_n,y_1,...,y_n) \quad \longrightarrow$ 
  \begin{align}
    &X, Y \sim \mathcal{N}(0,1)\\[10pt]
    &X \sim \mathcal{N}(0,1), \quad Y = a\,X + \sim \mathcal{N}(0,1)\\[10pt]
  \end{align}
\end{frame}

\begin{frame}
  Given two sets of observations $X_i$ and $Y_i$, Pearson's correlation
  coefficient $r$ is defined as

  \begin{equation*}
    r = \frac{\text{Cov}(X,Y)}{\sigma_X \sigma_Y} = 
    \frac{\sum_{i}(x_i-\bar{x})(y_i-\bar{y})}{\sqrt{\sum_{k}(x_k-\bar{x})^2}\sqrt{\sum_{l}(y_l-\bar{y})^2}}
  \end{equation*}\vspace{10pt}

  so that $r \in [-1,1]$.\\
  If $X, Y \sim \mathcal{N}$ independently (null hypothesis) the pdf for $r$ 
  has the exact distribution

  \begin{equation*}
    f(r) = \frac{\left(1-r^2\right)^{\frac{n-1}{2}}}{\sqrt{n-2}
    \;\;\text{B}\left(\frac{1}{2},\frac{n-2}{2}\right)}
  \end{equation*}
\end{frame}

\begin{frame}
  Given as set of observations $X_i$ we define the \textbf{ranks} of $X$ as 
  \begin{equation*}
    R(X) = \{ r(x_1), r(x_2), ..., r(x_n)\}, \quad r(x_i) \in \mathbb{N}
  \end{equation*}
  so that
  \begin{equation*}
    r(x_i) < r(x_j) \, \Leftrightarrow \, x_i \leq x_j\quad \text{and} \quad r(x_i)
    < \#X_i
  \end{equation*}
  
  \textbf{Sperman's $\rho$ correlation coefficient} between $X_i$ and $Y_i$ is the
  \textbf{linear correlation between the ranks $R(X)$ and $R(Y)$}

  \begin{equation*}
    \rho = \frac{\text{Cov}(R(X),R(Y))}{\sigma_{R(X)} \sigma_{R(Y)}}
  \end{equation*}\vspace{10pt}
\end{frame}

\begin{frame}
  \begin{equation*}
    Y = \sin(X) + \mathcal{N}(\mu, \sigma^2)
  \end{equation*}
  \begin{tabular}{c | c}
    $X$ & $Y$\\
    \hline
    0.5 & 2.0 \\
    4.1 & 2.4 \\
    2.7 & 1.5 \\
    2.2 & 3.0
  \end{tabular}
  $\longrightarrow$
  \begin{tabular}{c | c}
    $R(X)$ & $R(Y)$\\
    \hline
     1 & 2\\
     4 & 3\\
     3 & 1\\
     2 & 4
  \end{tabular}
  $\longrightarrow$
  \begin{tabular}{c | c l}
    $R(X)$ & $R(Y)$\\
    \hline
    1 & 2 & \scriptsize{$r_1$}\\
    2 & 4 & \scriptsize{$r_2$}\\
    3 & 1 & \scriptsize{$r_3$}\\
    4 & 3 & \scriptsize{$r_4$}\\
  \end{tabular}
  \begin{equation*}
    \xi_n(X,Y) = 1 - \frac{3 \sum_{i=1}^{n-1}{|r_{i+1}-r_i}|}{n^2 - 1}
  \end{equation*}
\end{frame}

\begin{frame}\footnotesize
  \textbf{Theorem 1.} If $Y$ is almost surely not a constant with law $\mu$.
  Then

  \begin{equation*}
    \xi_n(X,Y) \quad \xlongrightarrow[n \to \infty]{a.s.} \quad \xi(X,Y) = \frac{\int
    \Var{\E{\CF{\{Y \geq t\}}|X}}d\mu(t)}{\int \Var{\CF{\{Y \geq t\}}}d\mu(t)}
    \in [0,1]
  \end{equation*}

  with $\xi(X,Y)$ being $0$ if and only if $X$ and $Y$ are independent, and
  being $1$ if and only if there is a measurable function $f : \mathbb{R} \to
  \mathbb{R}$ such that $Y = f(X)$ almost surely.\\[20pt]

  \textbf{Theorem 2.} If $X$ and $Y$ are independent and $Y$ is
  continuous, then
  \begin{equation*}
    \sqrt{n} \; \xi_n(X,Y) \xlongrightarrow[n\to\infty]{distr} \mathcal{N}\left(0,\frac{2}{5}\right)
  \end{equation*}
  Importantly, this convergence is quite fast ($n \sim 20$).
\end{frame}

\begin{frame}\small
  \begin{enumerate}
    \item $\xi$ is \textbf{non symmetric}, but can be symmetrized with
      \begin{equation*}
        S_{\xi}(X,Y) = \max\left\{\xi(X,Y),\xi(Y,X)\right\}
      \end{equation*}
    \item \textbf{minimal assumptions}: $(X,Y)$ is i.i.d., $Y$ is non-constant
    \item \textbf{asymptotic boundedness:} ``it is clear that $\xi(X,Y) \in
      [0,1]$''\footnote{One has to be careful here} 
    \item \textbf{range}: in the absence of ties
      \begin{equation*}
        -\frac{1}{2} + O\left(\frac{1}{n}\right) \;\leq \; \xi_n(X,Y) \; \leq \;
        \frac{n-2}{n+1}
      \end{equation*}
    \item \textbf{There is a python package} with $\xi$ and new work build on
      top of it
  \end{enumerate}
\end{frame}

\begin{frame}
  \begin{equation*}
    Y = \sin(2\nu \pi X) + \cos(\nu \pi X) + \mathcal{N}(0,\sigma)
  \end{equation*}
  \begin{equation*}
    \begin{cases}
      &X \sim \mathcal{N}\\
      &Y \sim \mathcal{N}\\
      &Z = \cos^3(X) + \mathcal{N}\\
      &W = Y^4 + \mathcal{N}
    \end{cases}
  \end{equation*}

  \begin{equation*}
    \begin{cases}
      &X \sim \mathcal{N}\\
      &Y = \sin^2(X) + \mathcal{N}\\
      &Z = \cos^3(X) + \mathcal{N}\\
      &W = Y^4 + \mathcal{N}
    \end{cases}
  \end{equation*}

\end{frame}
\begin{frame}
  \begin{tabular}{ c | c | c }
    & \textbf{n} & \textbf{n}\\
    \hline
    $\xi_1$ & $X_n^{(1)}$ & $Y_n^{(1)}$ \\
    \hline
    $\xi_2$ & $X_n^{(2)}$ & $Y_n^{(2)}$ \\
    \hline
    ... & ... & ... \\
    \hline
    $\xi_{10^6}$ & $X_n^{(10^6)}$ & $Y_n^{(10^6)}$ \\
    \hline
  \end{tabular}
\end{frame}

\end{document}
